% Derek Sarmiento (scripts/01\_download\_mgn.py)
% Redactada por Miguel Duarte a partir del script, su self-check, el recibo y la
% corrida del 9 de noviembre de 2026; pendiente de revisión de Derek.
\section{Descarga del Marco Geoestadístico Nacional}\label{sec:descarga}

\textit{Responsable: Derek Sarmiento. Script: \f{scripts/01\_download\_mgn.py}.}

\subsection{Objetivo}

Obtener la capa municipal completa del MGN 2025 del DANE (1122 municipios), con los campos que usa el diccionario de datos de la Entrega~1 (\f{MPIO\_CDPMP}, \f{MPIO\_CNMBR}, \f{MPIO\_NAREA}, \f{DPTO\_CCDGO}), dejarla como un GeoPackage legible y acompañarla de un manifiesto que diga de dónde salieron los bytes y si la descarga es nacional y verificable.

\subsection{Fuente}

La fuente es el servicio WFS~2.0 público del DANE, capa \f{Serv\_CapasMGN\_2025:Municipio}, el acceso OGC que publica la ficha \f{8xdp-fsgm} de \f{datos.gov.co}. Se pide en GML~3.2 con \f{srsName=urn:ogc:def:crs:EPSG::4686}, en seis páginas de 200 registros (\f{startIndex} 0, 200, 400, 600, 800 y 1000). La página renderizada del geoportal y el \f{FeatureServer} REST no se usan: el cortafuegos del DANE los bloquea para clientes automáticos. La verificación TLS queda siempre activa.

\subsection{Entradas y salidas}

\begin{table}[H]
\centering\small
\caption{Contrato de entrada y salida de \f{01\_download\_mgn.py}.}
\begin{tabularx}{\textwidth}{@{}p{5.6cm} p{2.2cm} L@{}}
\toprule
Pieza & Rol & Condición \\
\midrule
\f{config.yaml} (\f{sources.mgn.wfs}) & entrada & servicio, \f{typeName}, \f{srsName}, plan de páginas y \f{expected\_feature\_count = 1122} \\
\f{outputs/mgn\_wfs\_source\_receipt.json} & entrada & por página: URL, bytes, \f{sha256}, \f{content\_sha256} y \f{server\_timestamp} \\
\f{data/raw/mgn2025/wfs/page\_XXXX.gml} & salida & respuestas del WFS guardadas byte a byte (fuera de Git) \\
\f{data/raw/mgn2025/wfs/municipio.gpkg} & salida & 1122 municipios, EPSG:4686, longitud/latitud (fuera de Git) \\
\f{outputs/mgn\_download\_manifest.json} & salida & \f{status = complete}, \f{claims\_national\_download = true}, SHA256 de la capa y coincidencia de cada página \\
\bottomrule
\end{tabularx}
\end{table}

\subsection{Método}

\begin{enumerate}
  \item \textbf{Plan oficial.} El recibo solo cuenta si su servicio, \f{typeName}, \f{srsName}, URLs, \f{startIndex} y \f{number\_matched} son los de \f{config.yaml}. Un recibo de otro tamaño, aunque sus bytes coincidan, no habilita una descarga nacional.
  \item \textbf{Caché o descarga.} Si las seis páginas locales coinciden con el recibo, no se contacta el servidor. Si no, se hace un intento por página y la corrida se detiene al primer fallo de red o TLS. Una respuesta que no coincide con el recibo se rechaza.
  \item \textbf{Coincidencia por bytes o por contenido.} El WFS escribe en el elemento raíz un atributo \f{timeStamp} con la hora de cada respuesta, así que dos descargas del mismo contenido nunca tienen el mismo SHA256. Por eso el recibo guarda además \f{content\_sha256}: el SHA256 de la página con el valor de ese único atributo vaciado. Una página vale si coincide por bytes o por contenido; los archivos descargados no se modifican.
  \item \textbf{Lectura del GML.} Los elementos vienen con alias en español y se mapean de forma explícita (\f{Código\_DANE\_concatenado\_departamento\_y\_municipio} $\rightarrow$ \f{MPIO\_CDPMP}, \f{Nombre\_municipio} $\rightarrow$ \f{MPIO\_CNMBR}, \f{Área\_departamento\_en\_kilómetros\_cuadrados} $\rightarrow$ \f{MPIO\_NAREA}, entre otros). Cada \f{posList} viene en latitud/longitud, como declara EPSG:4686, y se guarda en longitud/latitud.
  \item \textbf{Controles de la capa.} Exactamente 1122 registros, igual a \f{numberMatched}; códigos de 5 dígitos sin repetidos; un único \f{srsName} en todas las páginas.
  \item \textbf{Unidad del área.} El nombre del elemento dice kilómetros cuadrados; además se compara \f{MPIO\_NAREA} con el área geodésica WGS84 de cada contorno. Si la mediana del cociente no cae entre 0,5 y 2, no se declara la unidad y la descarga no queda como nacional.
  \item \textbf{Manifiesto honesto.} Si algo falla, el manifiesto actual queda en \f{blocked} o \f{invalid} con \f{claims\_national\_download = false}, y un manifiesto exitoso anterior se conserva aparte en \f{mgn\_download\_manifest.history.json}. Si la capa y los hash no cambiaron, la siguiente corrida solo la revalida.
\end{enumerate}

\subsection{Pruebas}

\f{python scripts/01\_download\_mgn.py --self-check} usa fixtures etiquetados y nunca cuenta como descarga del DANE. Resultado: \textbf{22/22 OK}. Cubre la elección del ZIP correcto en el HTML del geoportal, la lectura de la unidad desde metadatos, el rechazo de CRS ausente o distinto, de capas sin \f{MPIO\_NAREA} o ilegibles, la inversión de ejes del GML, que un recibo de un municipio no se anuncie como nacional, que un fallo TLS no deje el manifiesto anterior como actual, y los ocho casos del hash de contenido: coincidencia por bytes y por contenido, rechazo de un byte distinto fuera del \f{timeStamp}, rechazo sin \f{content\_sha256}, rechazo de dos \f{timeStamp} en la raíz, reconstrucción del SHA256 original y el \f{--backfill-receipt} aceptado y rechazado.

\subsection{Resultados}

Las respuestas originales se obtuvieron el 8 de noviembre de 2026 entre las 15:23 y las 15:24 (hora de Bogotá), en 6,4 a 8,5~s por página. En la noche del 8 de noviembre (21:38, hora de Bogotá) se volvieron a descargar las seis páginas: tienen el mismo tamaño y, con el \f{timeStamp} original reinsertado, reproducen exactamente el SHA256 del recibo. Es decir, son el mismo contenido; solo cambió la hora que escribe el servidor.

\begin{table}[H]
\centering\small
\caption{Páginas del recibo (\f{mgn\_wfs\_source\_receipt.json}).}
\begin{tabular}{@{}lrlll@{}}
\toprule
Página & Bytes & \f{sha256} & \f{content\_sha256} & \f{timeStamp} original \\
\midrule
\f{page\_0000.gml} & 26\,416\,139 & \f{5d600d25b928\ldots} & \f{ce7a0d0ca809\ldots} & 20:23:33Z \\
\f{page\_0200.gml} & 22\,174\,892 & \f{6900fa419817\ldots} & \f{09b5856d94dd\ldots} & 20:23:40Z \\
\f{page\_0400.gml} & 25\,069\,453 & \f{168772c3d7bd\ldots} & \f{89e051c376c9\ldots} & 20:23:47Z \\
\f{page\_0600.gml} & 29\,395\,276 & \f{d6b3ee0b193b\ldots} & \f{0c271831db52\ldots} & 20:23:54Z \\
\f{page\_0800.gml} & 21\,884\,852 & \f{2b8496abd8a2\ldots} & \f{31f8cfda05d0\ldots} & 20:24:02Z \\
\f{page\_1000.gml} & 27\,119\,583 & \f{2eadfebad833\ldots} & \f{11a39ac93152\ldots} & 20:24:09Z \\
\bottomrule
\end{tabular}
\end{table}

\begin{table}[H]
\centering\small
\caption{Capa resultante según \f{mgn\_download\_manifest.json}.}
\begin{tabularx}{\textwidth}{@{}p{5.6cm} L@{}}
\toprule
Indicador & Valor \\
\midrule
Estado & \f{complete}, \f{national\_complete = true}, \f{claims\_national\_download = true} \\
Registros & 1122 (igual a \f{numberMatched}), 0 geometrías nulas, todas \f{MultiPolygon} \\
Año de versión DANE & 954 municipios de 2024 y 168 de 2025 (no se filtra) \\
CRS & EPSG:4686 declarado por el WFS; coordenadas guardadas en longitud/latitud \\
Unidad de \f{MPIO\_NAREA} & km$^2$; cociente atributo/área geodésica: mediana 0,9994, mínimo 0,9984, máximo 1,0209 \\
Capa & \f{municipio.gpkg}, 100,6~MB \\
Coincidencia con el recibo & 6/6 páginas por contenido (descarga del 8 de noviembre, 21:38) \\
\bottomrule
\end{tabularx}
\end{table}

El SHA256 del GeoPackage cambia en cada escritura porque el formato guarda metadatos propios; el manifiesto registra el de la corrida actual y la preparación lo comprueba. Lo que se mantiene estable es el producto: a partir de estas páginas, la preparación produce siempre el mismo GeoJSONL (sección~\ref{sec:preparacion}).

\subsection{Limitaciones}

\begin{itemize}
  \item No se descargó el ZIP oficial del MGN ni su \f{.prj}, así que el CRS del archivo original no está probado (\f{original\_archive\_crs\_proven = false}). Se usa el \f{srsName} que declara el WFS; el almacenamiento REST del mismo servicio informa EPSG:3857.
  \item El elemento del área se llama \emph{Área\_departamento} aunque el objeto es un municipio; la unidad se confirma con el cociente geodésico, no con el nombre solo.
  \item Las páginas GML y el GeoPackage están fuera de Git por tamaño (unos 150~MB). Para reproducir la cadena basta con volver a descargarlas: el hash de contenido las valida contra el recibo.
\end{itemize}
